% ===================================================================
%  图表第 1 号 — 学校。— 中学。— 教室。
%
%  Ceci n'est PAS une transcription : c'est une traduction, faite en
%  2026, en chinois d'aujourd'hui. D'ou trois differences avec
%  texto/io et texto/fr, et trois seulement :
%
%   * pas de \nl ni de \cc — la traduction ne suit aucune ligne
%     imprimee, et les lignes de ce fichier ne sont que du confort ;
%   * pas de \begin{VUpage} — elle n'a ni feuillet ni folio, et la
%     page de lecture ne lui en affiche pas ;
%   * le gras se pose sur le TERME seul, ponctuation dehors.
%
%  Tout le reste est identique, et doit l'etre : les cles « %%K » sont
%  celles de texto/io, macro pour macro pour l'apparat, et les renvois
%  \textsuperscript{(n)} visent les MEMES objets de la planche — ce
%  sont eux qui ouvrent les gros plans.
%
%  L'ido est le texte de base ; le francais sert de controle, et
%  l'anglais et l'espagnol de calibrage.
%
%  UNE LANGUE SANS ESPACES ET SANS CASSE. Deux consequences pour la
%  page de lecture. La premiere : le renvoi se colle au terme, et le
%  blanc que uniformiser_renvois() pose entre les deux est le seul de
%  la ligne — il ne coupe rien, le chinois n'ayant pas de mot a
%  couper. La seconde : les regles de role d'apparat ne peuvent pas
%  s'appuyer sur la capitale. Le titre du tableau porte donc la suite
%  « 图表第 », que la prose n'ecrit jamais ; la note de l'editeur des
%  tableaux 3 et 4 dit « 第三幅和第四幅 » et ne s'y trompe pas.
% ===================================================================

%%K t01-apar-0 apar
\VUsaut{2.0mm}
\VUcentre{12.6pt}{120}{说明手册}
\VUsaut{3.2mm}
\VUcentre{8.4pt}{160}{附于}
\VUsaut{2.6mm}
\VUtitre{15.6pt}{0}{德尔马辅助挂图}
\VUsaut{3.4mm}
\VUornamento{9pt}
\VUsaut{5.0mm}
\VUcentre{13.2pt}{90}{第一组}
\VUsaut{3.0mm}
\VUfilet{16mm}
\VUsaut{5.6mm}

%%K t01-apar-1 apar
\VUtitre{15.0pt}{60}{图表第 1 号}
\VUsaut{1.4mm}
\VUtitre{11.4pt}{0}{学校。--- 中学。--- 教室。}
\VUsaut{2.2mm}
\VUfilet{20mm}
\VUsaut{6.4mm}

%%K t01-tit-1 sub
\VUpk{10.2pt}{40}{总的描述。}
\VUsaut{3.0mm}

%%K t01-01-1 p
1. --- 这幅图表画的是一所\VUgras{中学}里的\VUgras{教室}。教室里有八张
\VUgras{课桌} \textsuperscript{(18)} 和八条\VUgras{长凳} \textsuperscript{(32)}。每条长凳
上坐三个学生。\VUgras{老师} \textsuperscript{(7)} 坐在\VUgras{椅子} \textsuperscript{(8)} 上，在他
的\VUgras{讲台} \textsuperscript{(6)} 前。

%%K t01-02-1 p
2. --- 讲台左边立着一个带玻璃门的\VUgras{柜子} \textsuperscript{(15)}。右边是
\VUgras{黑板} \textsuperscript{(1)}，配着\VUgras{板擦} \textsuperscript{(2)} 和
\VUgras{粉笔} \textsuperscript{(3)}。

%%K t01-02-2 p
讲台上方挂着几幅\VUgras{挂图} \textsuperscript{(9, 11, 12)} 和一幅\VUgras{地图}
\textsuperscript{(10)}。

%%K t01-03-1 p
3. --- 并不是所有学生都在自己的\VUgras{座位} \textsuperscript{(17)} 上。约翰
\textsuperscript{(4)} 站在黑板前；他拿着板擦，正在擦去那些\VUgras{几何图形}。

%%K t01-03-2 p
彼得站在讲台旁；他正在收\VUgras{书面作业} \textsuperscript{(5)}，送给老师。

%%K t01-04-1 p
4. --- \VUgras{这位}老师是教\VUgras{现代语言}的。他教学生说、读、写外语
\VUgras{（德语、英语、西班牙语、意大利语、俄语}或\VUgras{法语）}。

%%K t01-05-1 p
5. --- 教室很大，也很明亮。尽头有一扇宽大的\VUgras{窗户}，带两扇
\VUgras{气窗} \textsuperscript{(78)}，还有一道\VUgras{玻璃门}，用来通风采光。

%%K t01-05-2 p
这间教室并不冷。屋子当中为了取暖，放着一个很好的\VUgras{炉子}
\textsuperscript{(46)}，连着它的\VUgras{烟囱} \textsuperscript{(45)}。

%%K t01-05-3 p
隔板上钉着\VUgras{衣帽钩} \textsuperscript{(58)}；学生的帽子和大衣挂在上面。

%%K t01-tit-2 sub
\VUsaut{5.2mm}
\VUpk{10.2pt}{40}{院子。}
\VUsaut{3.6mm}

%%K t01-06-1 p
6. --- \VUgras{时钟} \textsuperscript{(63)} 指着九点五分。

%%K t01-06-2 p
\VUgras{课间休息} \textsuperscript{(76)} 刚刚结束，课又开始了。休息是在
\VUgras{院子} \textsuperscript{(75)} 里过的。我们看见几个学生在玩。一位
\VUgras{学监} \textsuperscript{(74)} 在看着他们。

%%K t01-07-1 p
7. --- 院子里还看得见另外几个人：\VUgras{督学} \textsuperscript{(73)}、
\VUgras{校长} \textsuperscript{(71)} 和\VUgras{教务主任} \textsuperscript{(72)}。

%%K t01-07-2 p
督学视察各学校，校长管理这所学校，教务主任督查学生的功课。

%%K t01-07-3 p
第四个人是\VUgras{校工} \textsuperscript{(77)}。他腋下夹着一本册子，用来登记缺席
的人。

%%K t01-08-1 p
8. --- 院子尽头是\VUgras{体操器械} \textsuperscript{(70)} 和\VUgras{手工}
\textsuperscript{(69)} 作坊。

%%K t01-08-2 p
图表右边隐约看得见\VUgras{音乐}和\VUgras{图画}的\VUgras{教室}。

%%K t01-noto-1 noto
\VUnotes{143.2mm}{%
(*) 至于\textit{名字}，我们也建议按拉丁形式转写。只有这样才能免去数不清
的异写。何况在 \textit{Giovanni、Jean、Johann、Jan、Hans、John、Ivan}
之间又该怎么选呢？难道要另编一部名字词典不成？我们取拉丁语的主格，就说：
\VUgras{Ioannes, Ioanna; Iulius, Iulia; Iakobus; Andreas; Lukas.}
（《完全语法》）。%
}

%%K t01-tit-3 sub
\VUpk{10.2pt}{40}{详细的描述。}
\VUsaut{2.4mm}

%%K t01-tit-4 sub
\VUpk{10.2pt}{40}{好学生。}
\VUsaut{3.0mm}

%%K t01-09-1 p
9. --- 讲台右边第一条长凳上的学生是\VUgras{亨利}(*)、\VUgras{斯蒂芬}和
\VUgras{查理}。

%%K t01-09-2 p
亨利很专心，也很用功；他正在听老师讲。他桌上有一本\VUgras{练习本}
\textsuperscript{(28)}、一本\VUgras{书} \textsuperscript{(29)}、一本\VUgras{词典}
\textsuperscript{(30)} 和一个\VUgras{笔盒} \textsuperscript{(31)}。他的书有许多
\VUgras{页}；每一章有好几个\VUgras{段落}，每一\VUgras{行}有许多\VUgras{词}。

%%K t01-09-4 p
他的同桌斯蒂芬 \textsuperscript{(24)} 很勤快。他正把下次要交的功课记在本子上。
他的\VUgras{字}写得很漂亮。他的书合着。只看得见\VUgras{封面}
\textsuperscript{(27)}。他旁边放着他的\VUgras{圆规} \textsuperscript{(26)}。

%%K t01-09-5 p
查理 \textsuperscript{(23)} 手里拿着书；他正把它翻开，再温习一遍功课。跟每个学生
一样，他面前放着一个\VUgras{墨水瓶} \textsuperscript{(25)}。他很听话。

%%K t01-10-1 p
10. --- 下一张桌上坐着\VUgras{路易、安东}和\VUgras{保罗}。路易正在背他的
\VUgras{功课} \textsuperscript{(22)}，回答老师的\VUgras{问题}。安东
\textsuperscript{(21)} 很不专心；老师有时甚至罚他。保罗 \textsuperscript{(20)} 是个好
同伴，和气又肯帮忙。他很专心。

%%K t01-tit-5 sub
\VUsaut{4.2mm}
\VUpk{10.2pt}{40}{坏学生。}
\VUsaut{3.2mm}

%%K t01-11-1 p
11. --- 亨利后面坐着\VUgras{乔治} \textsuperscript{(38)}。他不是个坏孩子，可是
\VUgras{爱说话}、不专心，又坐不住。他的\VUgras{轻浮}和\VUgras{不听话}在
课上引起\VUgras{混乱}，他常常该受一顿严厉的\VUgras{告诫}。他的
\VUgras{草稿本} \textsuperscript{(35)} 很脏，他用\VUgras{钢笔} \textsuperscript{(34)} 在
上面弄出\VUgras{墨渍}，所以老是要用\VUgras{橡皮} \textsuperscript{(36)}；他的
\VUgras{书包} \textsuperscript{(33)} 破了，因为他太不小心。他把\VUgras{笔杆}
\textsuperscript{(37)} 带到现代语言教室来做什么呢？

%%K t01-11-2 p
他的同桌埃德蒙 \textsuperscript{(39)} 不喜欢他。埃德蒙确实有点儿懒，可是安静，
也坐得住。他不闲聊；只是他不听讲，宁愿去读\VUgras{读本}里的\VUgras{故事}。

%%K t01-11-3 p
这条长凳上第三个学生是\VUgras{吕多维克} \textsuperscript{(40)}。他很勤勉。他正在
一张白\VUgras{纸}上写字。他面前摊着\VUgras{吸墨纸} \textsuperscript{(41)}。

%%K t01-12-1 p
12. --- 老师左边第二条长凳上的学生年纪很小。头一个，小\VUgras{埃米尔}，
还不会好好写字；他带来他的\VUgras{石板} \textsuperscript{(42)}、他的\VUgras{石笔}
和他的\VUgras{海绵} \textsuperscript{(43)}。

%%K t01-12-2 p
他旁边是\VUgras{奥古斯特}和\VUgras{贝尔纳} \textsuperscript{(44)}。后者功课做得
好，可是\VUgras{样子}很邋遢。

%%K t01-13-1 p
13. --- 他们后面坐着三个\VUgras{懒鬼}。\VUgras{阿尔贝}不背功课。
\VUgras{雅各} \textsuperscript{(47)} 不听讲，倒睡着了。他的\VUgras{懒惰}给他招来
\VUgras{申斥}，甚至\VUgras{处罚}。最后，\VUgras{克里斯蒂安}
\textsuperscript{(48)} 太轻佻了。他\VUgras{品行}不好，也不肯改。

%%K t01-14-1 p
14. --- 他们的邻座却很守规矩。\VUgras{欧内斯特} \textsuperscript{(51)} 爱整齐、
守规矩；他总是用他的\VUgras{尺} \textsuperscript{(50)} 和他的\VUgras{铅笔}
\textsuperscript{(49)}。他是\VUgras{走读生}。

%%K t01-14-2 p
\VUgras{弗朗西斯} \textsuperscript{(52)} 是\VUgras{寄宿生}；他穿制服，在学校里吃
饭睡觉。他是个好\VUgras{朋友}。

%%K t01-14-3 p
莫里斯正用他的\VUgras{小刀} \textsuperscript{(56)} 削他的\VUgras{铅笔}；他很仔细。

%%K t01-14-4 p
他把书都带来了，他的\VUgras{语法书} \textsuperscript{(55)}、他的\VUgras{笔记本}
\textsuperscript{(54)}，还有一本\VUgras{便笺簿} \textsuperscript{(53)}。他的\VUgras{书袋}
\textsuperscript{(57)} 放在他身后的地上。

%%K t01-14-5 p
教室尽头的两张桌子坐的是\VUgras{好学生} \textsuperscript{(19)}。

%%K t01-tit-6 sub
\VUsaut{4.6mm}
\VUpk{10.2pt}{40}{图画和音乐。}
\VUsaut{3.4mm}

%%K t01-15-1 p
15. --- \VUgras{图画教室}里，墙上挂着漂亮的\VUgras{范本}，天花板上垂下一盏
带\VUgras{灯罩}的\VUgras{灯} \textsuperscript{(62)}。\VUgras{老师} \textsuperscript{(60)}
很有耐心；他给学生改\VUgras{画} \textsuperscript{(61)}，教他们照着\VUgras{胸像}
\textsuperscript{(59)} 写生。

%%K t01-16-1 p
16. --- \VUgras{音乐教室}看着很有意思。在那里学认\VUgras{乐谱}，唱写在
\VUgras{五线谱}上的\VUgras{音阶} \textsuperscript{(67)}（do、re、mi、fa、sol、
la、si\,=\,C. D. E. F. G. A. B.），认\VUgras{谱号}，唱动听的\VUgras{曲子}。
乐谱放在\VUgras{谱架} \textsuperscript{(65)} 上。一个学生在\VUgras{弹钢琴}
\textsuperscript{(64)}，\VUgras{音乐老师} \textsuperscript{(66)} 在\VUgras{打拍子}
\textsuperscript{(68)}。

%%K t01-17-1 p
17. --- 隐约还看得见\VUgras{手工} \textsuperscript{(69)} 作坊的一角，孩子们可以在
那里学刨木头、锉铁。并不是每所学校都有。

%%K t01-tit-7 sub
\VUsaut{5.0mm}
\VUpk{10.2pt}{40}{理科教室。}
\VUsaut{3.6mm}

%%K t01-18-1 p
18. --- 数学老师教学生\VUgras{几何、}\VUgras{算术、算法}和\VUgras{米制}。
在挂图上我们看得见主要的几何图形：\VUgras{圆} \textsuperscript{(a)}、
\VUgras{正方形} \textsuperscript{(b)}、\VUgras{三角形} \textsuperscript{(c)}、
\VUgras{直线} \textsuperscript{(d)}。

%%K t01-18-2 p
我们还看得见一道\VUgras{算式}；它不是\VUgras{加法}，不是\VUgras{减法}，
也不是\VUgras{乘法}：它是\VUgras{除法} \textsuperscript{(e)}。

%%K t01-19-1 p
19. --- 在米制挂图上我们看得见：一根\VUgras{米尺} \textsuperscript{(a)}、一架
\VUgras{天平} \textsuperscript{(b)} 和它的\VUgras{砝码}、一\VUgras{升}
\textsuperscript{(e)}、一\VUgras{十升} \textsuperscript{(f)}、一\VUgras{分升}和一
\VUgras{立方米} \textsuperscript{(d)}。

%%K t01-19-2 p
学生也学\VUgras{博物} \textsuperscript{(11)} ——动物学 \textsuperscript{(a)}、植物学
\textsuperscript{(b)}、地质学 \textsuperscript{(c)} ——和\VUgras{历史} \textsuperscript{(12)}。

%%K t01-19-3 p
柜子里放着\VUgras{物理} \textsuperscript{(15)} 和\VUgras{化学} \textsuperscript{(16)} 的
仪器。

%%K t01-19-4 p
柜子上面放着：一个\VUgras{球} \textsuperscript{(14)}、一个\VUgras{圆锥}和一个
\VUgras{地球仪} \textsuperscript{(13)}。

%%K t01-19-5 p
挂图上方挂着一幅\VUgras{地图}，画着世界的五大部分：\VUgras{欧洲}
\textsuperscript{(g)}、\VUgras{亚洲} \textsuperscript{(e)}、\VUgras{非洲} \textsuperscript{(f)}、
\VUgras{美洲} \textsuperscript{(ab)} 和\VUgras{大洋洲} \textsuperscript{(h)}，还有两大
洋——\VUgras{太平洋} \textsuperscript{(c)} 和\VUgras{大西洋} \textsuperscript{(d)}。
\VUsaut{6.0mm}
\VUfilet{22mm}
